\section{Группы автоморфизмов и неподвижные поля}

В этом разделе все поля считаются подполями некоторого фиксированного поля, если не сказано обратное. Для расширения $E / F$ будем писать
\[
    \Aut(E/F) \coloneqq \{\sigma \in \Aut(E) \mid \sigma(a) = a \text{ для всех } a \in F\}.
\]
То есть это автоморфизмы поля $E$, которые поточечно оставляют элементы $F$ на месте.

\begin{definition}
Пусть $G \leq \Aut(E)$ --- группа автоморфизмов поля $E$. \emph{Полем неподвижных элементов} группы $G$ называется
\[
    E^G \coloneqq \{x \in E \mid \sigma(x) = x \text{ для всех } \sigma \in G\}.
\]
\end{definition}

\begin{stmt}
$E^G$ является подполем поля $E$.
\begin{proof}
Если $x, y \in E^G$, то для любого $\sigma \in G$ имеем
\[
    \sigma(x+y)=\sigma(x)+\sigma(y)=x+y,
    \qquad
    \sigma(xy)=\sigma(x)\sigma(y)=xy.
\]
Кроме того, $\sigma(1)=1$, а если $x \neq 0$, то
\[
    \sigma(x^{-1}) = \sigma(x)^{-1} = x^{-1}.
\]
Значит, $E^G$ замкнуто относительно сложения, умножения и взятия обратных ненулевых элементов.
\end{proof}
\end{stmt}

\begin{remark}
Если $E / F$ --- расширение и $G \leq \Aut(E/F)$, то $F \subseteq E^G$. Действительно, каждый автоморфизм из $G$ по определению оставляет все элементы $F$ неподвижными.
\end{remark}

\begin{theorem}[теорема Артина о неподвижном поле]
Пусть $G \leq \Aut(E)$ --- конечная группа автоморфизмов поля $E$. Тогда расширение $E / E^G$ конечно и
\[
    [E : E^G] = |G|.
\]
В частности,
\[
    G = \Aut(E / E^G).
\]
\begin{proof}
Обозначим $K = E^G$.

Сначала доказывается нетривиальная оценка
\[
    [E : K] \leq |G|.
\]
Идея такая: если бы нашлось больше чем $|G|$ линейно независимых над $K$ элементов поля $E$, то, подставляя их во все автоморфизмы из $G$, получили бы однородную систему линейных уравнений. Из минимального нетривиального решения и независимости различных автоморфизмов как отображений $E \to E$ получается противоречие. Это стандартная лемма Артина о линейной независимости автоморфизмов поля.

Теперь так как $G$ состоит из автоморфизмов, фиксирующих $K$, имеем
\[
    G \leq \Aut(E/K).
\]
Поэтому
\[
    |G| \leq |\Aut(E/K)|.
\]
С другой стороны, число $K$-гомоморфизмов $E \to E$ не превосходит степени расширения:
\[
    |\Aut(E/K)| \leq |\Hom_K(E,E)| \leq [E:K].
\]
Собирая оценки в цепочку, получаем
\[
    [E:K] \leq |G| \leq |\Aut(E/K)| \leq [E:K].
\]
Следовательно, все неравенства являются равенствами, откуда
\[
    [E:E^G] = |G|,
    \qquad
    G = \Aut(E/E^G).
\]
\end{proof}
\end{theorem}

\begin{remark}
Теорема Артина говорит, что если мы начинаем не с расширения полей, а с конечной группы автоморфизмов $G$ поля $E$, то поле неподвижных элементов $E^G$ автоматически устроено так, что $G$ является полной группой автоморфизмов расширения $E / E^G$.
\end{remark}

\section{Нормальные, сепарабельные и галуа-расширения}

\begin{definition}
Пусть $E / F$ --- алгебраическое расширение. Оно называется \emph{нормальным}, если для любого $\alpha \in E$ минимальный многочлен элемента $\alpha$ над $F$ раскладывается в $E[x]$ на линейные множители.
\end{definition}

\begin{definition}
Пусть $E / F$ --- алгебраическое расширение. Оно называется \emph{сепарабельным}, если для любого $\alpha \in E$ минимальный многочлен элемента $\alpha$ над $F$ не имеет кратных корней.
\end{definition}

\begin{definition}
Конечное расширение $E / F$ называется \emph{расширением Галуа}, если оно нормально и сепарабельно.
В этом случае группа
\[
    \Gal(E/F) \coloneqq \Aut(E/F)
\]
называется \emph{группой Галуа} расширения $E / F$.
\end{definition}

\begin{examplelist}
\item $\mathbb{Q}(\sqrt{2}) / \mathbb{Q}$ --- расширение Галуа. Минимальный многочлен $x^2 - 2$ имеет два различных корня $\pm\sqrt{2}$, оба лежат в $\mathbb{Q}(\sqrt{2})$.

\item $\mathbb{Q}(\sqrt[3]{2}) / \mathbb{Q}$ не является нормальным. Минимальный многочлен $x^3 - 2$ имеет корни
\[
    \sqrt[3]{2}, \qquad \omega\sqrt[3]{2}, \qquad \omega^2\sqrt[3]{2},
\]
где $\omega$ --- первообразный кубический корень из единицы. В поле $\mathbb{Q}(\sqrt[3]{2})$ лежит только вещественный корень, поэтому многочлен не раскладывается в нём полностью.

\item Поле
\[
    \mathbb{Q}(\sqrt[3]{2}, \omega)
\]
является полем разложения многочлена $x^3 - 2$ над $\mathbb{Q}$.
\end{examplelist}

\begin{theorem}
Пусть $E / F$ --- конечное расширение. Тогда $E / F$ является расширением Галуа тогда и только тогда, когда $E$ является полем разложения некоторого сепарабельного многочлена $f \in F[x]$.
\begin{proof}
\textbf{($\Leftarrow$)} Пусть $E$ --- поле разложения сепарабельного многочлена $f \in F[x]$. Тогда $E/F$ конечно. Нормальность следует из того, что все корни $f$ уже лежат в $E$: любой $F$-гомоморфизм $E$ в алгебраическое замыкание переводит корни $f$ в корни $f$, а значит переводит $E$ в себя. Сепарабельность следует из сепарабельности $f$: элементы $E$ выражаются через корни сепарабельного многочлена.

\textbf{($\Rightarrow$)} Пусть $E/F$ --- конечное нормальное сепарабельное расширение. Возьмём конечный набор порождающих
\[
    E = F(\alpha_1, \ldots, \alpha_n).
\]
Пусть $f_i \in F[x]$ --- минимальный многочлен $\alpha_i$ над $F$. Так как расширение нормально, каждый $f_i$ раскладывается в $E[x]$ на линейные множители. Так как расширение сепарабельно, эти линейные множители попарно различны. Тогда многочлен
\[
    f(x) = \prod_{i=1}^n f_i(x)
\]
сепарабелен, раскладывается в $E[x]$, и поле $E$ порождено его корнями. Значит, $E$ является полем разложения сепарабельного многочлена $f$.
\end{proof}
\end{theorem}

\begin{cor}
Если $E/F$ --- расширение Галуа, то
\[
    |\Gal(E/F)| = [E:F].
\]
\begin{proof}
Пусть $G = \Gal(E/F)$. Всегда имеем $F \subseteq E^G$. По теореме Артина
\[
    [E:E^G] = |G|.
\]
Так как $E/F$ нормально и сепарабельно, число $F$-автоморфизмов $E$ максимально возможно, то есть равно $[E:F]$. Поэтому $E^G = F$, и
\[
    |G| = [E:E^G] = [E:F].
\]
\end{proof}
\end{cor}

\begin{remark}
Равенство $|\Aut(E/F)| = [E:F]$ --- характерный признак ситуации Галуа. В произвольном конечном расширении автоморфизмов может быть меньше степени расширения. Например, у $\mathbb{Q}(\sqrt[3]{2})/\mathbb{Q}$ нет нетривиального автоморфизма: образ $\sqrt[3]{2}$ обязан быть корнем $x^3-2$ внутри того же поля, а там лежит только $\sqrt[3]{2}$.
\end{remark}

\section{Основная теорема теории Галуа}

Пусть теперь $E/F$ --- конечное расширение Галуа и
\[
    G = \Gal(E/F).
\]
Главная идея: промежуточным полям между $F$ и $E$ соответствуют подгруппы группы $G$.

\begin{definition}
Если $H \leq G$, то ему сопоставляется поле неподвижных элементов
\[
    H \longmapsto E^H.
\]
Если $L$ --- промежуточное поле,
\[
    F \subseteq L \subseteq E,
\]
то ему сопоставляется подгруппа
\[
    L \longmapsto \Gal(E/L).
\]
\end{definition}

\begin{theorem}[основная теорема теории Галуа]
Пусть $E/F$ --- конечное расширение Галуа и $G = \Gal(E/F)$. Тогда отображения
\[
    H \longmapsto E^H,
    \qquad
    L \longmapsto \Gal(E/L)
\]
задают взаимно обратное соответствие между подгруппами $H \leq G$ и промежуточными полями $F \subseteq L \subseteq E$.

Это соответствие меняет включения на противоположные:
\[
    H_1 \leq H_2 \quad \Rightarrow \quad E^{H_2} \subseteq E^{H_1},
\]
\[
    L_1 \subseteq L_2 \quad \Rightarrow \quad \Gal(E/L_2) \leq \Gal(E/L_1).
\]
Кроме того,
\[
    [E:E^H] = |H|,
    \qquad
    [E^H:F] = \frac{|G|}{|H|},
\]
и
\[
    \Gal(E/E^H) = H,
    \qquad
    E^{\Gal(E/L)} = L.
\]
\begin{proof}
Разберём основные шаги.

\textbf{1. Включения меняются на противоположные.}
Если $H_1 \leq H_2$, то чем больше группа автоморфизмов, тем больше условий неподвижности накладывается на элемент. Поэтому
\[
    E^{H_2} \subseteq E^{H_1}.
\]
Если $L_1 \subseteq L_2$, то автоморфизм, фиксирующий $L_2$, тем более фиксирует $L_1$, значит
\[
    \Gal(E/L_2) \leq \Gal(E/L_1).
\]

\textbf{2. Переход от подгруппы к полю и обратно.}
Пусть $H \leq G$. Тогда по теореме Артина
\[
    [E:E^H] = |H|
    \quad \text{и} \quad
    \Aut(E/E^H) = H.
\]
Так как $E/F$ --- Галуа, $G$ имеет порядок $[E:F]$, поэтому
\[
    [E^H:F] = \frac{[E:F]}{[E:E^H]}
             = \frac{|G|}{|H|}.
\]

\textbf{3. Переход от промежуточного поля к подгруппе и обратно.}
Пусть $F \subseteq L \subseteq E$. Так как $E/F$ --- нормальное и сепарабельное расширение, расширение $E/L$ также нормально и сепарабельно. Значит, $E/L$ --- расширение Галуа. Тогда
\[
    |\Gal(E/L)| = [E:L].
\]
По теореме Артина для группы $\Gal(E/L)$ имеем
\[
    [E:E^{\Gal(E/L)}] = |\Gal(E/L)| = [E:L].
\]
Но очевидно $L \subseteq E^{\Gal(E/L)}$, потому что все автоморфизмы из $\Gal(E/L)$ фиксируют $L$. Два промежуточных поля дают одинаковую степень над $E$, значит они равны:
\[
    E^{\Gal(E/L)} = L.
\]
Тем самым обе конструкции действительно взаимно обратны.
\end{proof}
\end{theorem}

\begin{remark}
Удобно запомнить соответствие так: большая подгруппа фиксирует больше элементов, поэтому даёт меньшее поле. Поэтому диаграмма включений переворачивается.
\end{remark}

\section{Нормальные подгруппы и фактор-группы}

Пусть снова $E/F$ --- конечное расширение Галуа, $G = \Gal(E/F)$, и $H \leq G$. Обозначим
\[
    L = E^H.
\]

\begin{stmt}
Для любого $\sigma \in G$ выполнено
\[
    E^{\sigma H \sigma^{-1}} = \sigma(E^H).
\]
\begin{proof}
Пусть $x \in E^H$. Докажем, что $\sigma(x)$ неподвижен относительно всех элементов $\sigma H \sigma^{-1}$.
Возьмём $h \in H$. Тогда
\[
    (\sigma h \sigma^{-1})(\sigma(x)) = \sigma(h(x)) = \sigma(x).
\]
Значит, $\sigma(E^H) \subseteq E^{\sigma H \sigma^{-1}}$. Обратное включение получается применением того же рассуждения к $\sigma^{-1}$.
\end{proof}
\end{stmt}

\begin{theorem}
Пусть $E/F$ --- конечное расширение Галуа, $G = \Gal(E/F)$, $H \leq G$ и $L = E^H$. Тогда следующие условия равносильны:
\begin{enumerate}
    \item $H \trianglelefteq G$.
    \item $L/F$ --- расширение Галуа.
\end{enumerate}
Если эти условия выполнены, то
\[
    \Gal(L/F) \cong G/H.
\]
\begin{proof}
Рассмотрим отображение ограничения
\[
    \rho\colon G \to \Aut(L/F),
    \qquad
    \rho(\sigma) = \sigma|_L.
\]
Оно корректно тогда, когда каждый $\sigma \in G$ переводит $L$ в себя. По предыдущему утверждению
\[
    \sigma(L) = \sigma(E^H) = E^{\sigma H \sigma^{-1}}.
\]
Следовательно, $\sigma(L)=L$ для всех $\sigma \in G$ ровно тогда, когда
\[
    \sigma H \sigma^{-1}=H \quad \forall \sigma \in G,
\]
то есть когда $H$ нормальна в $G$.

Если $H \trianglelefteq G$, то $L$ сохраняется всеми автоморфизмами из $G$, и расширение $L/F$ получается нормальным. Сепарабельность наследуется от расширения $E/F$, значит $L/F$ является расширением Галуа.

Найдём ядро отображения ограничения:
\[
    \ker \rho = \{\sigma \in G \mid \sigma|_L = \mathrm{id}_L\} = \Gal(E/L) = H.
\]
По основной теореме о гомоморфизме групп
\[
    G/H \cong \rho(G) \leq \Gal(L/F).
\]
Но
\[
    |G/H| = \frac{|G|}{|H|} = [L:F] = |\Gal(L/F)|,
\]
поэтому включение является равенством:
\[
    \Gal(L/F) \cong G/H.
\]
\end{proof}
\end{theorem}

\begin{remark}
В основной теореме теории Галуа нормальные подгруппы соответствуют именно тем промежуточным полям, которые сами дают расширения Галуа над исходным полем $F$.
\end{remark}
